\renewcommand{\thetable}{3.A.6}
\begin{tablalafrox}{Supuestos de la oferta. Fuente: elaboración propia a partir de los requerimientos de Distribuidora Puelche y del Anexo 2.2 del Subdocumento 2 (LafroX, 2026a) y sección «Arbitraje de tensiones operacionales y comerciales» del Subdocumento 2 (LafroX, 2026c).}%
  {tab:3-A-6}%
  {F{0.045} F{0.245} Y F{0.200} F{0.130} F{0.145}}%
  {\cab{ID} & \cab{Supuesto adoptado} & \cab{Fundamento} & \cab{Impacto si resulta falso} & \cab{Validación} & \cab{Origen}}
  S-01 & OTIF considera cumplida solo la entrega del 100 \% de ítems y unidades dentro de la fecha y ventana pactadas. & Las áreas calculan el OTIF de forma distinta. La definición adoptada en el SD2 permite medir el cumplimiento con un criterio común (LafroX, 2026a, Anexo 2.2, S-01). & Cambia la fórmula de RF-11.03 y la línea base de R18-04; se recalculan las metas. & Gerencia Comercial y de Operaciones, mes 1. & SD2, Anexo 2.2, S-01 \\
  S-02 & La unidad de trazabilidad sanitaria es el lote del proveedor; el SSCC identifica la unidad logística asociada. & El retiro sanitario exige identificar los productos y clientes vinculados al lote del proveedor; identificar únicamente la unidad logística no permite resolver esa consulta. & Se redefine la captura en recepción (RF-01.02, RF-01.07) y la consulta de retiro. & Jefa de Calidad, mes 1. & SD2, Anexo 2.2, S-02 \\
  S-03 & Ante local cerrado, el conductor registra el intento y el sistema reagenda a la próxima ventana disponible; no se entrega a terceros. Si el cliente continúa ausente, la mercadería retorna al centro de distribución con trazabilidad del evento. & Los locales cerrados requieren un tratamiento que conserve la trazabilidad de la entrega fallida y de la mercadería, conforme al flujo adoptado en el SD2 (LafroX, 2026a, Anexo 2.2, S-03). & Cambia RNG-05 y el flujo de reentrega de M6 y M4. & Gerencia de Operaciones, levantamiento de reglas. & SD2, Anexo 2.2, S-03 \\
  S-04 & La excursión térmica se parametriza por tipo de producto. Una excursión menor y transitoria genera una alerta preventiva en cabina; una excursión crítica y sostenida hace que el sistema bloquee preventivamente el lote y notifique a Calidad, que decide liberar, bloquear o rechazar. & El SD2 adopta una respuesta térmica graduada para conciliar continuidad operacional y control sanitario, manteniendo en Calidad la decisión final sobre el producto (LafroX, 2026a, Anexo 2.2, S-04). & Si la autoridad exige invalidación automática, cambia RNG-04 y el tratamiento en ruta. & Jefa de Calidad, con los parámetros por producto, mes 2. & SD2, Anexo 2.2, S-04; SD2, arbitraje 2 \\
  S-05 & Cada transportista confirma conductor y vehículo antes del despacho; el conductor real se autentica al iniciar la ruta. & Los conductores externos rotan sin aviso. Confirmar la asignación y autenticar al conductor real permite atribuir los registros a la persona, vehículo y viaje correspondientes. & Sin confirmación previa se incorpora la vinculación en el andén (RF-14.05). & Acuerdo operacional con cada una de las 10 empresas transportistas. & SD2, Anexo 2.2, S-05 \\
  S-06 & Se mantiene el efectivo y se reduce mediante medios alternativos voluntarios; cada camión realiza rendición digital individual y toda diferencia exige causal tipificada y responsable antes del cierre. & El canal tradicional debe conservar el pago en efectivo. La rendición individual permite investigar las diferencias sin imponer un cambio de medio de pago. & Si no se logra reducir el uso de efectivo mediante medios alternativos voluntarios, se mantiene la rendición digital individual y se ajusta la capacidad de conciliación al volumen efectivo de cobros, sin exigir al canal tradicional abandonar el pago en efectivo. & Gerente de Finanzas, mes 1. & SD2, Anexo 2.2, S-06;  \\
  S-07 & El costo de servir se calcula por actividad, entrega y cliente, con distancia, tiempos, volumen, devoluciones y retornables. Los clientes no rentables se segmentan para revisar frecuencia y pedido mínimo; no se eliminan automáticamente. & El prorrateo por zona con un criterio de 2016 no refleja el costo de cada atención. El SD2 adopta un cálculo basado en las actividades y recursos utilizados (LafroX, 2026a, Anexo 2.2, S-07). & Otra metodología cambia RF-11.02 y los datos que se capturan en E1. & Gerente de Finanzas, antes del mes 13. & SD2, Anexo 2.2, S-07 \\
  S-08 & El stock se reserva al confirmar el pedido en el servidor central, por orden cronológico; sin conexión el stock es indicativo. & La captura sin conexión impide garantizar reservas comunes en tiempo real. La confirmación central evita que varios dispositivos comprometan el mismo stock. & Una reserva local cambia RNG-01 y la resolución de conflictos. & Gerencia Comercial, levantamiento de reglas. & SD2, Anexo 2.2, S-08 \\
  S-09 & Rige el precio acordado por el preventista al tomar el pedido y registrado en este, conforme a las condiciones comerciales autorizadas. Los cambios posteriores de la lista de precios no modifican el importe acordado con el cliente. & La política adoptada en el Subdocumento 2 conserva las condiciones comerciales acordadas al tomar el pedido y evita que una actualización posterior altere el importe pactado (LafroX, 2026a, Anexo 2.2, S-09). & Otra política exige revisar RNG-08 y la integración con el ERP, que debe conservar el precio registrado en el pedido para la facturación. & Gerencia Comercial, mes 2. & SD2, Anexo 2.2, S-09 \\
  S-10 & Canastillos y pallets se controlan por saldo de cliente y transportista, sin serialización unitaria. & El SD2 adopta el control por saldos para registrar entregas, devoluciones y diferencias de envases por cliente y transportista, sin exigir identificación individual (LafroX, 2026a, Anexo 2.2, S-10). & La serialización exige otra captura en M8 y otro hardware. & Gerencia de Operaciones, mes 2. & SD2, Anexo 2.2, S-10 \\
  S-11 & El maestro interno de Puelche prevalece y conserva equivalencias e historial de cambios de códigos externos. Los códigos sin equivalencia quedan en una bandeja de excepciones y no crean una segunda verdad de producto. & Las equivalencias e historial permiten relacionar códigos externos con el producto interno, evitando registros duplicados y cambios de identidad por diferencias de codificación. & Otro maestro rector cambia RNG-11 y RF-12.03. & Jefe de TI y Gerencia Comercial, mes 2. & SD2, Anexo 2.2, S-11 \\
  S-12 & La devolución se registra en la entrega con SKU, cantidad y motivo; la solución transmite el evento y la evidencia, y el ERP emite la nota de crédito cuando corresponde. & La devolución ocurre en terreno, pero el ERP conserva la emisión tributaria. Separar la captura del evento de su tratamiento tributario respeta esa responsabilidad. & Se mantiene el ERP como único emisor en cualquier alternativa. & Gerente de Finanzas, mes 2. & SD2, Anexo 2.2, S-12;  \\
  S-13 & No se instalan cámaras en cabina. Se integra la telemetría existente para gestionar vehículo, ruta y carga, con controles de privacidad que consideran que la ubicación del vehículo puede revelar la del conductor. El uso del GPS para control de jornada queda sujeto a acuerdo previo con el sindicato. & La objeción sindical exige distinguir la telemetría operacional del control de jornada. El SD2 adopta esa separación y condiciona este último uso a un acuerdo previo (LafroX, 2026a, Anexo 2.2, S-13). & Si no se alcanza el acuerdo, se mantiene la telemetría operacional con controles de privacidad y no se utiliza GPS para control de jornada. & Jefe de TI y Gerencia de Operaciones, durante el levantamiento, para definir finalidades y controles de privacidad; sindicato, antes de habilitar el uso del GPS para control de jornada. & SD2, Anexo 2.2, S-13;  \\
  S-14 & El WMS de 2013 se reemplaza por la solución común, con migración y reversión. & El WMS de 2013 tiene soporte discontinuado y coexiste con registros fragmentados. El SD2 adopta su reemplazo, preservando los datos y la posibilidad de reversión (LafroX, 2026a, Anexo 2.2, S-14). & Si se conserva, se agrega una integración y cambia la migración. & Jefa de Bodega y Jefe de TI, mes 3. & SD2, Anexo 2.2, S-14; Anexo 2.3 \\
  S-15 & La prueba principal es la firma en pantalla; en su defecto, fotografía, nombre del receptor o confirmación QR vinculada a la GDE. La alternativa conserva fecha, ubicación e identidad disponible y se articula con el acuse tributario. & La evidencia de entrega debe capturarse en terreno y admitir alternativas cuando no sea posible firmar, conservando su vinculación con la guía y el acuse tributario. & Otra evidencia cambia RF-06.02 y RF-06.11. & Gerente de Finanzas, mes 2. & SD2, Anexo 2.2, S-15 \\
  S-16 & El conocimiento del planificador se captura en la Etapa 1 y se parametriza en el motor de rutas antes de su jubilación. LafroX propone talleres en los meses 1 a 3. & La planificación depende del conocimiento de una persona próxima a jubilarse. Transferir y probar sus reglas tempranamente reduce esa dependencia operacional. & Si el planificador se retira antes, la ruta se valida con su ayudante y con datos históricos. & Gerencia de Operaciones y planificador, mes 1. & SD2, Anexo 2.2, S-16; SD2, arbitraje 5 \\
  S-17 & El plan se gobierna por meses relativos del Art. 17°, y el contrato se inicia en un mes que no ubica los meses 16 ni 21 en septiembre o diciembre (febrero, marzo, mayo, julio, agosto, octubre, noviembre o diciembre). & Los pasos a producción de los meses contractuales 16 y 21 deben compatibilizarse con la prohibición de realizarlos en septiembre y diciembre. & Un inicio en enero, abril, junio o septiembre genera una incompatibilidad entre los meses contractuales de paso a producción y las ventanas del caso. El asunto se plantea mediante V-17 durante el período oficial de consultas y su resolución requiere una respuesta formal del mandante. & Contraparte Técnica, al firmar el contrato. & BA, Art. 17°; Caso, sección 13.3, condición 2 \\
  S-18 & Las interfaces del ERP existen y pueden usarse sin modificar el ERP. & La solución debe integrarse con el ERP sin modificarlo. La existencia y utilización de interfaces se adopta como hipótesis pendiente de verificar durante el levantamiento. & Sin interfaz disponible se usa intercambio de archivos controlado y se reprograma la integración. & Jefe de TI, mes 1. &  \\
  S-19 & El CLIENTE puede acceder a los datos de telemetría del proveedor actual. & Existe un servicio de telemetría para la flota propia. Su integración requiere acceso a los datos, cuyas condiciones deben confirmarse con el proveedor. & LafroX propone un mecanismo alternativo con igual cobertura de posición, velocidad y kilometraje de los 42 camiones propios, sujeto a autorización del CLIENTE; la cobertura de RF-14.01 no se reduce. & Jefe de TI y proveedor de telemetría, mes 2. &  \\
  S-20 & El hardware de terreno lo compra el CLIENTE con la especificación de LafroX, antes de cada ola. & La compra del hardware corresponde al CLIENTE y su disponibilidad condiciona cada ola. La especificación de LafroX permite coordinar la adquisición con el despliegue. & La ola afectada se posterga dentro de la holgura disponible y se aplica un plan de recuperación para mantener los hitos contractuales. Si la demora supera esa holgura, se registra y escala el impacto con medidas de mitigación, sin asumir una modificación automática de los hitos. & Gerencia de Operaciones, según calendario de olas. & ) \\
  S-21 & La promesa de entrega es de 24 horas para clientes urbanos con pedidos ingresados antes de las 14:00, y de 48 horas para clientes rurales, periféricos o abastecidos por cross-docking. & El SD2 resuelve la tensión entre Comercial y Operaciones mediante una promesa diferenciada por ubicación y abastecimiento, con un corte horario para pedidos urbanos (LafroX, 2026c, sección «Arbitraje de tensiones operacionales y comerciales», arbitraje 1). & Si la capacidad resulta insuficiente, se revisan la planificación y los recursos para sostener la promesa y el corte definidos en el SD2. & Gerencia Comercial y de Operaciones, mes 2. & SD2, arbitraje 1 \\
  S-22 & Las instalaciones operacionales se confirman al inicio; la solución se parametriza por sitio (RF-02.01). & Las entrevistas indican cinco instalaciones y la volumetría del caso señala seis. Confirmarlas al inicio permite definir los sitios y las olas del despliegue. & Cambia el número de sitios con borde operacional y de olas. & Gerencia de Operaciones, mes 1. & SD2, Anexo 2.2, E-06 (discrepancia entre cinco y seis instalaciones) \\
  S-23 & La dotación de reparto comprende 84 conductores y peonetas propios y aproximadamente 160 conductores de transportistas: unas 244 personas, conforme al SD2. Esta cifra dimensiona la población a incorporar y capacitar; la asignación de equipos de reparto por camión se declara en S-30. & El SD2 identifica 84 conductores y peonetas propios y aproximadamente 160 conductores externos. Estas poblaciones determinan el alcance de incorporación y capacitación (LafroX, 2026c, sección «Resumen Ejecutivo del problema»). & Cambia la población a incorporar y capacitar; se revisa el efecto sobre dispositivos compartidos y simultaneidad. & Gerencia de Operaciones, mes 2. & SD2, «Resumen Ejecutivo del problema» y Anexo 2.3 \\
  S-24 & Se diseña para 14 horas sin señal y 24 horas sin enlace, aunque los cortes relatados sean menores. Es un mínimo obligatorio del caso y se trata como restricción (Anexo 3.E, R-02 y R-03). & El caso exige mínimos de autonomía de 14 horas en terreno y 24 horas en los centros de distribución; los cortes relatados no reducen esa obligación. & No aplica: es el mínimo del caso. & No requiere validación. &  \\
  S-25 & La reposición usa la historia de ventas disponible en el ERP y las promociones comprometidas. & La reposición requiere considerar la demanda histórica y las promociones comprometidas. La disponibilidad y calidad de la historia del ERP deben comprobarse durante el levantamiento. & Con menos historia se amplían las excepciones revisables de RF-10.01. & Jefe de abastecimiento, mes 3. & Caso, sección 4.1 \\
  S-26 & La pérdida anual de envases no supera el 7 \% del parque, verificada tras 12 meses de operación, conforme al SD2. La línea base es el 14 \% estimado del parque al año. El control por saldo de cliente y transportista permite registrar movimientos y detectar diferencias. & El SD2 fija una pérdida anual máxima del 7 \% del parque frente al 14 \% estimado actual. Su comprobación requiere controlar movimientos y medir un año de operación (LafroX, 2026c, Tabla 2.1). & Si la medición arroja otra línea base, se actualiza el diagnóstico y se revisan las medidas de control; se conserva la meta de pérdida anual no superior al 7 \% del parque. & Gerencia de Operaciones y Gerente de Finanzas, mes 2, para validar la medición. & SD2, Tabla 2.1 y Anexo 2.2, S-10; Caso, secciones 4.8 y 7.2, y cap. 18, resultado 12 \\
  S-27 & LafroX compromete una ocupación mínima del 60 \% de la capacidad volumétrica útil de cada camión al salir a reparto y una ocupación promedio superior al 68 \% actual. El mínimo se aplica también a las rutas rurales, sin excepciones por baja densidad. La planificación debe cumplir simultáneamente este umbral, las restricciones de peso y condición térmica y la promesa de entrega de S-21. & El caso exige aumentar la ocupación y comprometer un mínimo por camión. LafroX adopta el 60 \% como meta propia, compatible con las restricciones de carga y la promesa del SD2. & Si la validación operacional muestra que la planificación propuesta no permite cumplir simultáneamente la ocupación mínima y la promesa de entrega, se revisan la asignación de vehículos, la consolidación de cargas y las rutas, sin autorizar despachos bajo el umbral ni ampliar los plazos comprometidos. & Gerencia de Operaciones, con participación de la Gerencia Comercial, mes 2, mediante análisis de cargas, capacidades y plazos por ruta, incluidas las rurales. & Caso, secciones 4.4 y 7.2, y capítulo 18, resultado 14; SD2, arbitraje 1; RNG-14 y S-21. \\
  S-28 & Las empresas transportistas aceptan que se instale un termógrafo en sus 10 camiones con equipo de frío, y el termógrafo queda funcionando antes de que esas rutas partan con el nuevo proceso de reparto. El CLIENTE compra el equipo y LafroX lo especifica e integra. & Distribuidora Puelche no controla los vehículos de terceros; el alcance exige registro continuo de temperatura en los vehículos con frío. & Esos 10 camiones quedan sin registro propio y habría que exigir el dato por contrato al transportista. & Gerencia de Operaciones y transportistas, mes 6. & La empresa no controla los vehículos de terceros (Bases Técnicas del . 5), y el caso exige registro continuo de temperatura en los vehículos (cap. 9, p. 18). \\
  S-29 & La discrepancia del conteo cíclico actual es de 2,3 \% del valor contado. Conforme al SD2, las diferencias de inventario se investigan y los ajustes se justifican; ese porcentaje es línea base del problema y no una tolerancia automática de migración. Documentos y lotes deben coincidir sin diferencias. & La diferencia entre inventario contable y existencias físicas evidencia la necesidad de conciliar los datos antes de migrarlos. El SD2 exige investigar las diferencias y justificar los ajustes; la discrepancia histórica no autoriza aceptar errores en la migración (LafroX, 2026c, Tabla 2.1). & Si el conteo arroja otra discrepancia, se actualiza el diagnóstico y se revisa el saneamiento; se mantiene la investigación y justificación de cada diferencia. & Contraparte Técnica, antes de cada corte. & SD2, Tabla 2.1 \\
  S-30 & Los equipos de reparto (terminal del conductor, impresora de cabina y terminal de pago) se asignan por camión y no por conductor. La flota base comprende 42 camiones propios y 54 de transportistas, 96 en total, conforme al SD2. & Los conductores de los transportistas rotan sin aviso y Distribuidora Puelche no controla sus vehículos ni a sus conductores. & Una asignación personal requiere determinar quién usa cada tipo de equipo dentro de las 84 personas propias y aproximadamente 160 conductores externos de S-23; esas 244 personas incluyen peonetas y no equivalen a 244 conductores. & Levantamiento con Operaciones y transportistas, meses 1 a 3. & SD2, «Resumen Ejecutivo del problema» y Anexo 2.3;  \\
  S-31 & La flota de camiones crece en la misma proporción que los viajes mensuales. & La proyección operacional contempla los viajes a tres años, pero no especifica el crecimiento de la flota. & Cambia la compra por crecimiento de los equipos de reparto; la compra inicial no varía. & Revisión anual de capacidad. & El caso proyecta los viajes a tres años, pero no la flota (Bases Técnicas del . 24). \\
  S-32 & Los terminales de preventa y de bodega crecen en la misma proporción que la dotación de preventistas y del personal de los centros de distribución. & La proyección operacional contempla esas dotaciones a tres años. & Cambia la compra por crecimiento de esos terminales. & Revisión anual de capacidad. & El caso proyecta esas dotaciones a tres años (Bases Técnicas del . 24). \\
  S-33 & La carga de los camiones la hace la misma cuadrilla nocturna que prepara los pedidos, sin personal adicional con terminal. & La operación descrita integra la preparación y la carga en el proceso del turno de noche. & Habría que sumar terminales para el personal de carga. & Levantamiento en bodega, mes 1. & El caso describe la preparación y la carga como un mismo proceso del turno de noche (Bases Técnicas del . 8). \\
  S-34 & Las líneas de congelado son proporcionales a la participación del congelado en los productos de frío y en el área de las cámaras. & Los antecedentes operacionales informan los productos de frío y la superficie de cada cámara, pero no las líneas por zona. & Faltan o sobran terminales aptos para $-22\,^{\circ}\mathrm{C}$; el total de terminales no cambia. & Medición de líneas por zona en la marcha blanca. & El caso informa los productos de frío y la superficie de cada cámara, pero no las líneas por zona (Bases Técnicas del . 4–5). \\
  S-35 & En cada cross-docking trabajan dos personas a la vez durante la desconsolidación. & Los antecedentes operacionales indican personal reducido en los cross-docking, pero no precisan su cantidad. & Cambia la cantidad de terminales de los cross-docking. & Levantamiento en las plataformas, Etapa 1. & El caso indica personal reducido, pero no cuántas personas (Bases Técnicas del . 6). \\
  S-36 & Un punto de acceso Wi-Fi cubre unos 500\,m\textsuperscript{2} de bodega con racks, y la cobertura dentro de las cámaras exige mayor densidad. & Las cámaras actualmente no tienen señal; la cantidad definitiva de puntos de acceso se determina mediante el estudio de cobertura. & Cambia la cantidad de puntos de acceso y de switches de acceso. & Estudio de cobertura, Etapa 1. & Las cámaras hoy no tienen señal (Bases Técnicas del . 10), y la cantidad definitiva sale del estudio de cobertura exigido (RT-03.24; Bases Técnicas del . 26). \\
  S-37 & La cámara de frío de Concepción tiene una superficie proporcional a la del sitio respecto de Talca, unos 450\,m\textsuperscript{2}. & Los antecedentes operacionales indican que Concepción dispone de refrigerado, pero no precisan la superficie de su cámara. & Cambia la cantidad de puntos de temperatura y de puntos de acceso de Concepción. & Levantamiento en Concepción, mes 1. & El caso informa que Concepción tiene refrigerado, pero no el tamaño de su cámara (Bases Técnicas del . 5). \\
  S-38 & La flota de camiones con equipo de frío no crece en el horizonte de cotización. & Los antecedentes operacionales identifican 28 camiones con frío, pero no proyectan su crecimiento. & Hay que comprar termógrafos adicionales. & Revisión anual de capacidad. & El caso informa 28 camiones con frío, pero no proyecta su crecimiento (Bases Técnicas del . 5). \\
  S-39 & La preparación nocturna de Concepción ocupa la mitad de las personas que la de Talca, es decir, 60 preparadores a la vez. & Los antecedentes operacionales identifican 120 preparadores nocturnos en Talca y describen Concepción como un centro de menor escala, con la mitad de la superficie, pero no precisan su dotación. & Faltan o sobran terminales de bodega en Concepción, y cambia su capacidad local. & Levantamiento en Concepción, mes 1. & El caso informa 120 preparadores nocturnos en Talca (Bases Técnicas del . 14) y describe Concepción como un centro de menor escala, con la mitad de la superficie (cap. 2, p. 5, y cap. 3, p. 6), pero no da su dotación. \\
  S-40 & Una eventual séptima instalación se equipa con la misma tipología de Concepción y no se cotiza hasta que el CLIENTE defina el sitio. & La proyección operacional contempla siete instalaciones y un posible centro de distribución en Los Lagos hacia 2030, sin precisar ubicación ni superficie. & Compra adicional de equipamiento cuando se defina el sitio. & Confirmación del sitio por el CLIENTE. & El caso proyecta siete instalaciones y menciona un posible centro de distribución en Los Lagos hacia 2030, sin ubicación ni superficie (Bases Técnicas del . 22, y cap. 14, p. 24). \\
  S-41 & Cada una de las 184 personas de administración, comercial y soporte usa un computador propio, que se incorpora a la gestión central y a CrowdStrike. & Los antecedentes operacionales identifican 184 personas de administración, comercial y soporte en Talca y Concepción, pero no precisan cuántos computadores tienen. Cada equipo que entra a la plataforma debe estar gestionado y protegido (\textit{Bases Técnicas Transversales}, 2026, cap. 8, RT-08.09, p. 19). & Cambia la cantidad de equipos gestionados y protegidos. & Inventario de TI del CLIENTE, mes 1. & El caso informa 184 personas de administración, comercial y soporte en Talca y Concepción (Bases Técnicas del . 5), pero no cuántos computadores tienen. Cada equipo que entra a la plataforma debe estar gestionado y protegido (RT-08.09; Bases Técnicas Transversales, cap. 8, p. 19). \\
\end{tablalafrox}
